\section{Horton型地表流と飽和地表流の定義を分ける。}

\subsection{平成25年度 問29}
\label{q:h25-reg-applied-29}

\noindent\textbf{原典：} sources/H25/17\_applied\_science.pdf、PDF物理10頁。

\textbf{公式正答：} 1（single）。

\paragraph{問題の要旨} Horton型地表流と飽和地表流の定義を分ける。

\Needspace{4\baselineskip}\paragraph{確認した条件・記述（独自要約）}\mbox{}\par

\noindent\textbf{原典確認：} S7\_049の26原典画像で全五肢・否定・添字を確認。\par

\noindent\textbf{対象：} Horton型地表流と飽和地表流の定義を分ける。\par

\paragraph{問われる核心} Horton型地表流と飽和地表流の定義を分ける。

\paragraph{必要知識} 浸透能超過と飽和超過、Darcy則、TとKと固有透過度の区別。

\paragraph{考え方・正答を選ぶ理由（査読待ち初稿）} 地表流の発生を浸透能と飽和状態の二つで整理する。降雨強度が浸透能を上回るものはHorton型であり1が誤り。Darcy流束q=−K dh/dlと河川への流入を区別する。K=\allowbreak{}kρg/μ、一定KでT=\allowbreak{}Kbである。

\Needspace{4\baselineskip}\paragraph{選択肢の検討（査読待ち初稿）}\mbox{}\par

\noindent\textbf{1：不適切} 降雨が浸透能を超える流れを飽和地表流とする。これは浸透能超過型、すなわちHorton型。飽和型は土壌が飽和して追加の水を受け入れられなくなる場合である。\par

\noindent\textbf{2：通常の条件で適切} 透水量係数T=\allowbreak{}Kb。水平流を扱いKを厚さ方向に一定とする場合の式。一般には厚さ方向のKの積分となる。\par

\noindent\textbf{3：通常の条件で適切} 透水係数は媒体固有だけではない。Kは媒体の固有透過度に加え流体の密度と粘性に依存する。\par

\noindent\textbf{4：通常の条件で適切} 得水河川は地下水から水を受ける。地下水の流入により河川流量が増す区間をいう。\par

\noindent\textbf{5：通常の条件で適切} 絶対間隙率と有効間隙率の区別。総間隙体積の割合と、流れに有効な連結間隙の割合を分ける。\par

\paragraph{再利用する考え方} 測定対象と分離対象、定義の分母、絶対量と相対校正、全称と成立条件を分けて判断する。

\paragraph{共通理論} \hyperref[theory:hydro:flow]{第\ref*{theory:hydro:flow}節}

\par 参考：\cite{analysis49_runoff_ref,analysis49_darcy_ref}。

\par 難易度 3（校準初稿）、優先度未評価。

\subsection{令和元年度再試験 問29}
\label{q:r01-ret-applied-29}

\noindent\textbf{原典：} sources/R01\_retest/17\_applied\_science.pdf、PDF物理10頁。

\textbf{公式正答：} 5（single）。

\par\textbf{独立解答と公式正答に相違あり：} 独立解答 1|5。公式5、独立初回1|5を保持。原典1にのみ／だけの限定を追加して公式へ合わせない。Darcy流速と実流速の意味及び扇状地の地質・水頭条件も確認が必要。

\paragraph{問題の要旨} Darcy則・異方性・井戸水位と自噴条件を比較する。

\Needspace{4\baselineskip}\paragraph{確認した条件・記述（独自要約）}\mbox{}\par

\noindent\textbf{原典確認：} S7\_049の26原典画像で全五肢・否定・添字を確認。\par

\noindent\textbf{対象：} Darcy則・異方性・井戸水位と自噴条件を比較する。\par

\paragraph{問われる核心} Darcy則・異方性・井戸水位と自噴条件を比較する。

\paragraph{必要知識} 公式5は扇状地の地下構造と被圧水頭の代表的な説明である。1の原典は被圧の場合に成立するという文で、被圧の場合にのみという限定はない。Darcy則の一般性を指摘するだけでは1が誤りとは言えない。2～4は異方性、透水量係数の逆比例、水頭による流向から排除できるが、1と5の単一選択の整合は査読事項として残す。

\paragraph{考え方・正答を選ぶ理由（査読待ち初稿）} 公式5は扇状地の地下構造と被圧水頭の代表的な説明である。1の原典は被圧の場合に成立するという文で、被圧の場合にのみという限定はない。Darcy則の一般性を指摘するだけでは1が誤りとは言えない。2～4は異方性、透水量係数の逆比例、水頭による流向から排除できるが、1と5の単一選択の整合は査読事項として残す。

\Needspace{4\baselineskip}\paragraph{選択肢の検討（査読待ち初稿）}\mbox{}\par

\noindent\textbf{1：条件付き・解釈留保} Darcy則は被圧帯水層で成立する。適切な層流・連続体近似のもとでは被圧帯水層にも成立する。原典に被圧だけ、のみという語はない。非被圧でも使えることは、この肢を直接否定しない。\par

\noindent\textbf{2：不適切} 水平・鉛直の透水性の違いは稀。堆積層の配列や亀裂方向による異方性は一般的であり、等方性は近似条件である。\par

\noindent\textbf{3：不適切} 一定揚水の水位低下は厚さに比例。均一な被圧帯水層の定常井戸モデルではs=\allowbreak{}Q ln(R/r)/(2πKb)。厚さbへの依存は反比例である。\par

\noindent\textbf{4：不適切} 漏水は常に下向き。流れの向きは上下帯水層の水頭差で決まり、上向きも生じ得る。\par

\noindent\textbf{5：条件付き・解釈留保} 扇状地下流の被圧層と自噴。透水層を難透水層が覆い、水頭が地表より高い場合に自噴する。扇端等でその構造が成立する代表例であり、すべての扇状地に必ず起こるわけではない。\par

\paragraph{再利用する考え方} 測定対象と分離対象、定義の分母、絶対量と相対校正、全称と成立条件を分けて判断する。

\paragraph{共通理論} \hyperref[theory:hydro:flow]{第\ref*{theory:hydro:flow}節}

\par 参考：\cite{analysis49_darcy_ref,inorganic_hydro_ref}。

\par 難易度 4（校準初稿）、優先度未評価。

\section{全地球水に占める氷・地下水を含む淡水の割合を選ぶ。}

\subsection{平成24年度 問35}
\label{q:h24-reg-applied-35}

\noindent\textbf{原典：} sources/H24/17\_applied\_science.pdf、PDF物理10頁。

\textbf{公式正答：} 1（single）。

\par\textbf{検討未確定：} 水量は推計・対象範囲に依存する代表値。3\%を厳密な物理定数とはしない。

\paragraph{問題の要旨} 全地球水に占める氷・地下水を含む淡水の割合を選ぶ。

\Needspace{4\baselineskip}\paragraph{確認した条件・記述（独自要約）}\mbox{}\par

\noindent\textbf{原典確認：} S7\_049の26原典画像で全五肢・否定・添字を確認。\par

\noindent\textbf{対象：} 全地球水に占める氷・地下水を含む淡水の割合を選ぶ。\par

\paragraph{問われる核心} 全地球水に占める氷・地下水を含む淡水の割合を選ぶ。

\paragraph{必要知識} 分母は河川水だけでも利用可能な淡水だけでもなく、海水を含む地球上の水全体である。USGSの2010年資料では淡水約2.5\%なので、提示値で最も近い3\%を選ぶ。氷河・永久雪や地下水を淡水の内訳から除かない。

\paragraph{考え方・正答を選ぶ理由（査読待ち初稿）} 分母は河川水だけでも利用可能な淡水だけでもなく、海水を含む地球上の水全体である。USGSの2010年資料では淡水約2.5\%なので、提示値で最も近い3\%を選ぶ。氷河・永久雪や地下水を淡水の内訳から除かない。

\Needspace{4\baselineskip}\paragraph{選択肢の検討（査読待ち初稿）}\mbox{}\par

\noindent\textbf{1：通常の条件で適切} 3\%。全水量を分母とする代表推計約2.5\%に最も近い。\par

\noindent\textbf{2：不適切} 5\%。淡水総量の代表割合約2.5\%より大きい。\par

\noindent\textbf{3：不適切} 7\%。淡水総量の代表割合約2.5\%より大きい。\par

\noindent\textbf{4：不適切} 9\%。淡水総量の代表割合約2.5\%より大きい。\par

\noindent\textbf{5：不適切} 11\%。淡水総量の代表割合約2.5\%より大きい。\par

\paragraph{再利用する考え方} 測定対象と分離対象、定義の分母、絶対量と相対校正、全称と成立条件を分けて判断する。

\paragraph{共通理論} \hyperref[theory:hydro:global]{第\ref*{theory:hydro:global}節}

\par 参考：\cite{analysis49_fresh_ref}。

\par 難易度 4（校準初稿）、優先度未評価。

\section{比貯留率と貯留係数の次元と厚さ依存を確認する。}

\subsection{令和7年度 問19}
\label{q:r07-reg-applied-19}

\noindent\textbf{原典：} sources/R07/17\_applied\_science.pdf、PDF物理7頁。

\textbf{公式正答：} 4（single）。

\paragraph{問題の要旨} 比貯留率と貯留係数の次元と厚さ依存を確認する。

\Needspace{4\baselineskip}\paragraph{確認した条件・記述（独自要約）}\mbox{}\par

\noindent\textbf{原典確認：} S7\_049の26原典画像で全五肢・否定・添字を確認。\par

\noindent\textbf{対象：} 比貯留率と貯留係数の次元と厚さ依存を確認する。\par

\paragraph{問われる核心} 比貯留率と貯留係数の次元と厚さ依存を確認する。

\paragraph{必要知識} 比貯留率Ssは単位体積・単位水頭の放出量で次元は長さの逆数、被圧帯水層の貯留係数S=\allowbreak{}Ss bは単位面積・単位水頭で無次元である。4は体積と面積の定義を入れ替え、厚さの掛け方も逆にしている。T=\allowbreak{}Kbの次元は長さの二乗／時間。次元の比較で名称暗記の取り違えを防げる。

\paragraph{考え方・正答を選ぶ理由（査読待ち初稿）} 比貯留率Ssは単位体積・単位水頭の放出量で次元は長さの逆数、被圧帯水層の貯留係数S=\allowbreak{}Ss bは単位面積・単位水頭で無次元である。4は体積と面積の定義を入れ替え、厚さの掛け方も逆にしている。T=\allowbreak{}Kbの次元は長さの二乗／時間。次元の比較で名称暗記の取り違えを防げる。

\Needspace{4\baselineskip}\paragraph{選択肢の検討（査読待ち初稿）}\mbox{}\par

\noindent\textbf{1：通常の条件で適切} 空気侵入の負圧水頭。飽和から吸引を増して空気が入る閾値。正の吸引量として表す別の符号規約と区別する。\par

\noindent\textbf{2：通常の条件で適切} 透水係数は媒体固有だけではない。K=\allowbreak{}kρg/μで、流体の密度と粘性の影響も含む。\par

\noindent\textbf{3：通常の条件で適切} 含水比は水質量／乾燥固相質量。通常百分率で表す。体積含水率や湿潤総質量を分母とする割合とは異なる。\par

\noindent\textbf{4：不適切} 単位体積の貯留係数を厚さ倍して比貯留率。単位体積・単位水頭の量は比貯留率Ss。厚さを掛けて無次元貯留係数S=\allowbreak{}Ss bとなり、名前と関係が逆。\par

\noindent\textbf{5：通常の条件で適切} 透水量係数はKと厚さの積。一定Kの水平帯水層でT=\allowbreak{}Kb。SI次元は面積／時間である。\par

\paragraph{再利用する考え方} 測定対象と分離対象、定義の分母、絶対量と相対校正、全称と成立条件を分けて判断する。

\paragraph{共通理論} \hyperref[theory:hydro:storage]{第\ref*{theory:hydro:storage}節}

\par 参考：\cite{analysis49_storage_ref,analysis49_darcy_ref,inorganic_hydro_ref}。

\par 難易度 3（校準初稿）、優先度未評価。

\section{地下水の水質・水頭・含水量の定義を比較する。}

\subsection{平成26年度 問24}
\label{q:h26-reg-applied-24}

\noindent\textbf{原典：} sources/H26/17\_applied\_science.pdf、PDF物理10頁。

\textbf{公式正答：} 2（single）。

\paragraph{問題の要旨} 地下水の水質・水頭・含水量の定義を比較する。

\Needspace{4\baselineskip}\paragraph{確認した条件・記述（独自要約）}\mbox{}\par

\noindent\textbf{判断対象：} 地下水の水質・水頭・含水量の定義を比較する。\par

\noindent\textbf{原典確認：} 22原典頁画像で五候補・否定・化学式を照合。\par

\paragraph{問われる核心} 地下水の水質・水頭・含水量の定義を比較する。

\paragraph{必要知識} 地下水水質は起源の水だけでなく岩石との溶解・沈殿、酸化還元、生物・人為作用等を反映するので2。

\paragraph{考え方・正答を選ぶ理由（査読待ち初稿）} 地下水水質は起源の水だけでなく岩石との溶解・沈殿、酸化還元、生物・人為作用等を反映するので2。全水頭は位置と圧力の和、動水勾配は無次元。θは全体積を分母にし、飽和時θ=n、e=\allowbreak{}n/(1−n)。定義の名前だけでなく分母・次元を点検する。

\Needspace{4\baselineskip}\paragraph{選択肢の検討（査読待ち初稿）}\mbox{}\par

\noindent\textbf{1：選ばない候補} 帯水層に不飽和を含み被圧に自由水面。飽和し水を伝える帯水層と不飽和帯、自由水面の不圧を取り違え。\par

\noindent\textbf{2：選ぶ候補} 水質は降水/地表水/地層/環境に依存。起源と水岩反応・環境条件が水質を変える。\par

\noindent\textbf{3：選ばない候補} 全水頭は位置−圧力、勾配を拡散率。低速地下水のh=\allowbreak{}z+p/(ρg)は和で、差/距離は動水勾配。\par

\noindent\textbf{4：選ばない候補} 体積含水率は水/間隙、飽和で間隙比。θ=Vw/V全体、飽和では間隙率n。間隙比eではない。\par

\noindent\textbf{5：選ばない候補} 基底流は低水期除外、都市で直接流出減。基底流は低水期にも地下水等が支える流量。都市流出の断定も逆。\par

\paragraph{再利用する考え方} 原子番号と周期・族、物性と反応性、酸強度と危険性、生成と除去、定義の分母を別々に確かめる。

\paragraph{共通理論} \hyperref[theory:hydro:definitions]{第\ref*{theory:hydro:definitions}節}

\par 参考：\cite{inorganic_waterchem,inorganic_hydro_ref}。

\paragraph{関連問題} \hyperref[q:h27-reg-applied-23]{H27-REG-APPLIED-23}

\par 難易度 2（校準初稿）、優先度未評価。

\subsection{平成27年度 問23}
\label{q:h27-reg-applied-23}

\noindent\textbf{原典：} sources/H27/17\_applied\_science.pdf、PDF物理9頁。

\textbf{公式正答：} 3（single）。

\par\textbf{検討未確定：} 大気/河川平均滞留時間は資料・期間・推計の定義に依存。原典8日とUSGS約10日の違いを保持。

\paragraph{問題の要旨} 含水比の分母と陸水の移動過程を比較する。

\Needspace{4\baselineskip}\paragraph{確認した条件・記述（独自要約）}\mbox{}\par

\noindent\textbf{判断対象：} 含水比の分母と陸水の移動過程を比較する。\par

\noindent\textbf{原典確認：} 22原典頁画像で五候補・否定・化学式を照合。\par

\paragraph{問われる核心} 含水比の分母と陸水の移動過程を比較する。

\paragraph{必要知識} 通常の土質含水比は水質量/乾燥固体質量で、湿潤土全体で割った3は別の質量分率。

\paragraph{考え方・正答を選ぶ理由（査読待ち初稿）} 通常の土質含水比は水質量/乾燥固体質量で、湿潤土全体で割った3は別の質量分率。同じ試料でも値が異なる。失水/得水河川は交換の向き、Horton流は浸透能力を越える雨、蒸発散は気化の合算。大気滞留約8日の値を全推計の厳密定数とせず、比較関係と数値の留保を分ける。

\Needspace{4\baselineskip}\paragraph{選択肢の検討（査読待ち初稿）}\mbox{}\par

\noindent\textbf{1：選ばない候補} 大気平均滞留8日で河川より短い。約8日は推計条件依存。USGSには大気約10日・河川16–26日例もある。\par

\noindent\textbf{2：選ばない候補} 失水河川は地下へ浸透。河川から帯水層へ水が失われる区間。\par

\noindent\textbf{3：選ぶ候補} 含水比は水/湿潤土全体。通常土質含水比w=\allowbreak{}Mw/M乾燥であり湿潤全体ではない。\par

\noindent\textbf{4：選ばない候補} 浸透能より降雨強度大ならHorton流。浸透能力を超え余剰が地表へ流れる。\par

\noindent\textbf{5：選ばない候補} 蒸散と蒸発を合わせ蒸発散。植物経由と非植物面の気化を合算。\par

\paragraph{再利用する考え方} 原子番号と周期・族、物性と反応性、酸強度と危険性、生成と除去、定義の分母を別々に確かめる。

\paragraph{共通理論} \hyperref[theory:hydro:definitions]{第\ref*{theory:hydro:definitions}節}

\par 参考：\cite{inorganic_hydro_ref,inorganic_waterbudget}。

\paragraph{関連問題} \hyperref[q:h26-reg-applied-24]{H26-REG-APPLIED-24}

\par 難易度 2（校準初稿）、優先度未評価。
